\chapter{Advanced Database Objects and Optimization}
\label{chap:advanced}

\section{Indexes}

Indexes are created on columns frequently used in \texttt{WHERE}, \texttt{JOIN}, and \texttt{ORDER BY} clauses to reduce full-table scan costs on large tables.

\begin{lstlisting}[style=sqlstyle, caption={Index Definitions}]
-- Speed up investor lookups in transaction queries
CREATE INDEX idx_trans_investor  ON Transactions(Investor_ID);
CREATE INDEX idx_trans_date      ON Transactions(Date);
CREATE INDEX idx_trans_ticker    ON Transactions(Ticker);

-- Composite index for portfolio history range queries
CREATE INDEX idx_port_investor_date 
    ON Portfolio_History(Investor_ID, Date);

-- Speed up BehaviorSignals queries by investor
CREATE INDEX idx_signal_investor 
    ON BehaviorSignals(InvestorID, ObservationDate);
\end{lstlisting}

% TODO: Add EXPLAIN output screenshot showing index usage
% \begin{figure}[H]
%     \centering
%     \includegraphics[width=0.9\textwidth]{figures/explain_output.png}
%     \caption{EXPLAIN output before and after index creation}
% \end{figure}

\section{Views}

Three views are defined to support common reporting queries and dashboard data feeds.

\subsection{Investor Selling Summary}

\begin{lstlisting}[style=sqlstyle, caption={View: vw\_investor\_sell\_summary}]
CREATE VIEW vw_investor_sell_summary AS
SELECT 
    t.Investor_ID,
    i.Investor_Type,
    COUNT(CASE WHEN t.Action = 'SELL' THEN 1 END)   AS Total_Sells,
    COUNT(CASE WHEN t.Action = 'BUY'  THEN 1 END)   AS Total_Buys,
    COUNT(*)                                          AS Total_Trades,
    ROUND(COUNT(CASE WHEN t.Action = 'SELL' THEN 1 END) 
          / COUNT(*) * 100, 2)                        AS Sell_Ratio_Pct,
    AVG(bs.PanicScore)                                AS Avg_PanicScore
FROM Transactions t
JOIN Investors i ON t.Investor_ID = i.Investor_ID
LEFT JOIN BehaviorSignals bs ON t.Investor_ID = bs.InvestorID
GROUP BY t.Investor_ID, i.Investor_Type;
\end{lstlisting}

\subsection{Abnormal Sell Periods}

\begin{lstlisting}[style=sqlstyle, caption={View: vw\_abnormal\_sell\_periods}]
CREATE VIEW vw_abnormal_sell_periods AS
SELECT 
    Date,
    Ticker,
    COUNT(*)                            AS Sell_Count,
    SUM(Quantity * Price)               AS Total_Sell_Value
FROM Transactions
WHERE Action = 'SELL'
  AND Reason LIKE 'PANIC%'
GROUP BY Date, Ticker
HAVING Sell_Count > (
    SELECT AVG(cnt) + 2 * STDDEV(cnt)
    FROM (
        SELECT Date, COUNT(*) AS cnt
        FROM Transactions
        WHERE Action = 'SELL'
        GROUP BY Date
    ) sub
);
\end{lstlisting}

\subsection{Warning Dashboard}

\begin{lstlisting}[style=sqlstyle, caption={View: vw\_warning\_dashboard}]
CREATE VIEW vw_warning_dashboard AS
SELECT 
    w.WarningID,
    w.InvestorID,
    i.Name,
    i.Investor_Type,
    w.WarningDate,
    w.PanicLevel,
    w.Confidence,
    w.KeySignals,
    bs.DrawdownLevel,
    bs.SellSpike,
    bs.PanicScore
FROM Warnings w
JOIN Investors i        ON w.InvestorID = i.Investor_ID
JOIN BehaviorSignals bs ON w.InvestorID = bs.InvestorID
    AND w.WarningDate = bs.ObservationDate
ORDER BY w.PanicScore DESC, w.WarningDate DESC;
\end{lstlisting}

\section{User-Defined Functions}

Three UDFs encapsulate the core behavioral metric calculations.

\subsection{Rolling Drawdown}

\begin{lstlisting}[style=sqlstyle, caption={UDF: fn\_rolling\_drawdown}]
DELIMITER $$
CREATE FUNCTION fn_rolling_drawdown(
    p_investor_id VARCHAR(20),
    p_end_date    DATE,
    p_days        INT
) RETURNS FLOAT DETERMINISTIC
BEGIN
    DECLARE v_peak  FLOAT;
    DECLARE v_final FLOAT;
    
    SELECT MAX(Total_Asset) INTO v_peak
    FROM Portfolio_History
    WHERE Investor_ID = p_investor_id
      AND Date BETWEEN DATE_SUB(p_end_date, INTERVAL p_days DAY) 
                   AND p_end_date;
    
    SELECT Total_Asset INTO v_final
    FROM Portfolio_History
    WHERE Investor_ID = p_investor_id
      AND Date = p_end_date
    LIMIT 1;
    
    IF v_peak IS NULL OR v_peak = 0 THEN RETURN 0.0; END IF;
    RETURN ROUND((v_peak - v_final) / v_peak, 4);
END$$
DELIMITER ;
\end{lstlisting}

\subsection{Sell Intensity}

\begin{lstlisting}[style=sqlstyle, caption={UDF: fn\_sell\_intensity}]
DELIMITER $$
CREATE FUNCTION fn_sell_intensity(
    p_investor_id VARCHAR(20),
    p_start_date  DATE,
    p_end_date    DATE
) RETURNS FLOAT DETERMINISTIC
BEGIN
    DECLARE v_panic_sells INT;
    DECLARE v_total_trades INT;
    
    SELECT COUNT(*) INTO v_panic_sells
    FROM Transactions
    WHERE Investor_ID = p_investor_id
      AND Action = 'SELL'
      AND Reason LIKE 'PANIC%'
      AND Date BETWEEN p_start_date AND p_end_date;
    
    SELECT COUNT(*) INTO v_total_trades
    FROM Transactions
    WHERE Investor_ID = p_investor_id
      AND Date BETWEEN p_start_date AND p_end_date;
    
    IF v_total_trades = 0 THEN RETURN 0.0; END IF;
    RETURN ROUND(v_panic_sells / v_total_trades, 4);
END$$
DELIMITER ;
\end{lstlisting}

\subsection{Panic Score Aggregation}

\begin{lstlisting}[style=sqlstyle, caption={UDF: fn\_panic\_score\_agg}]
DELIMITER $$
CREATE FUNCTION fn_panic_score_agg(
    p_drawdown        FLOAT,
    p_sell_intensity  FLOAT,
    p_loss_sensitivity FLOAT
) RETURNS FLOAT DETERMINISTIC
BEGIN
    -- Weighted sum: drawdown 40%, sell intensity 40%, loss sensitivity 20%
    DECLARE v_score FLOAT;
    SET v_score = (p_drawdown * 0.4) 
                + (p_sell_intensity * 0.4) 
                + (p_loss_sensitivity * 0.2);
    -- Clamp to [0, 1]
    IF v_score > 1.0 THEN SET v_score = 1.0; END IF;
    IF v_score < 0.0 THEN SET v_score = 0.0; END IF;
    RETURN ROUND(v_score, 4);
END$$
DELIMITER ;
\end{lstlisting}

\section{Trigger}

The trigger automatically inserts a warning record when a new BehaviorSignal is inserted with a Panic Score exceeding 0.65.

\begin{lstlisting}[style=sqlstyle, caption={Trigger: trg\_panic\_warning}]
DELIMITER $$
CREATE TRIGGER trg_panic_warning
AFTER INSERT ON BehaviorSignals
FOR EACH ROW
BEGIN
    IF NEW.PanicScore >= 0.65 THEN
        INSERT INTO Warnings (
            InvestorID, WarningDate, PanicLevel, Confidence, KeySignals
        )
        VALUES (
            NEW.InvestorID,
            NEW.ObservationDate,
            CASE 
                WHEN NEW.PanicScore >= 0.80 THEN 'High'
                WHEN NEW.PanicScore >= 0.65 THEN 'Medium'
                ELSE 'Low'
            END,
            ROUND(NEW.PanicScore, 4),
            JSON_OBJECT(
                'drawdown',        NEW.DrawdownLevel,
                'sell_spike',      NEW.SellSpike,
                'loss_sensitivity', NEW.LossSensitivity,
                'panic_score',     NEW.PanicScore
            )
        );
    END IF;
END$$
DELIMITER ;
\end{lstlisting}

\section{Stored Procedures}

\subsection{Refresh Behavior Signals}

\begin{lstlisting}[style=sqlstyle, caption={SP: sp\_refresh\_behavior\_signals}]
DELIMITER $$
CREATE PROCEDURE sp_refresh_behavior_signals(
    IN p_investor_id VARCHAR(20),
    IN p_obs_date    DATE
)
BEGIN
    DECLARE v_drawdown       FLOAT;
    DECLARE v_sell_intensity FLOAT;
    DECLARE v_loss_sens      FLOAT;
    DECLARE v_panic_score    FLOAT;
    DECLARE v_start_date     DATE;
    
    SET v_start_date = DATE_SUB(p_obs_date, INTERVAL 30 DAY);
    
    -- Calculate each component using UDFs
    SET v_drawdown       = fn_rolling_drawdown(p_investor_id, p_obs_date, 30);
    SET v_sell_intensity = fn_sell_intensity(p_investor_id, v_start_date, p_obs_date);
    SET v_loss_sens      = fn_sell_intensity(p_investor_id, v_start_date, p_obs_date);
    SET v_panic_score    = fn_panic_score_agg(v_drawdown, v_sell_intensity, v_loss_sens);
    
    -- Insert or update BehaviorSignals (trigger fires automatically)
    INSERT INTO BehaviorSignals 
        (InvestorID, ObservationDate, DrawdownLevel, SellSpike, LossSensitivity, PanicScore)
    VALUES 
        (p_investor_id, p_obs_date, v_drawdown, v_sell_intensity, v_loss_sens, v_panic_score)
    ON DUPLICATE KEY UPDATE
        DrawdownLevel   = v_drawdown,
        SellSpike       = v_sell_intensity,
        LossSensitivity = v_loss_sens,
        PanicScore      = v_panic_score;
END$$
DELIMITER ;
\end{lstlisting}

\subsection{Simulate New Transactions (Demo)}

\begin{lstlisting}[style=sqlstyle, caption={SP: sp\_simulate\_new\_transactions}]
DELIMITER $$
CREATE PROCEDURE sp_simulate_new_transactions(
    IN p_investor_id VARCHAR(20),
    IN p_n_trades    INT
)
BEGIN
    DECLARE i INT DEFAULT 0;
    DECLARE v_action VARCHAR(10);
    DECLARE v_ticker VARCHAR(10);
    DECLARE v_price  FLOAT;
    
    WHILE i < p_n_trades DO
        -- Random SELL-heavy simulation to trigger panic detection
        SET v_action = IF(RAND() < 0.7, 'SELL', 'BUY');
        SET v_ticker = ELT(FLOOR(1 + RAND() * 9), 
                           'VIC','VHM','HPG','FPT','VCB',
                           'SSI','VPB','BCM','MSN');
        SET v_price  = 20000 + RAND() * 80000;
        
        INSERT INTO Transactions 
            (Trans_ID, Investor_ID, Date, Ticker, Action, Price, Quantity, Reason)
        VALUES (
            UUID(), p_investor_id, CURDATE(), v_ticker, v_action,
            ROUND(v_price, 2), FLOOR(100 + RAND() * 900), 
            CONCAT('PANIC_', v_action)
        );
        SET i = i + 1;
    END WHILE;
    
    -- Recalculate signals → trigger fires → warning auto-inserted
    CALL sp_refresh_behavior_signals(p_investor_id, CURDATE());
END$$
DELIMITER ;
\end{lstlisting}
